\newglossaryentry{term-agent}{
    name={Интеллектуальный программный агент},
    description={программная система, которая по входным данным самостоятельно планирует и выполняет последовательность действий для достижения поставленной цели, при необходимости обращаясь к внешним инструментам и пользователю}
}
\newglossaryentry{term-graph}{
    name={Граф состояний агента},
    description={ориентированный граф, вершины которого соответствуют узлам-обработчикам, а дуги~--- условным переходам между ними; определяет порядок и условия выполнения шагов обработки запроса}
}
\newglossaryentry{term-essential}{
    name={Существенные условия договора},
    description={условия, без согласования которых договор по российскому гражданскому праву не считается заключённым: предмет договора, условия, прямо названные таковыми в законе для соответствующего типа сделки, и условия, относительно которых по заявлению одной из сторон должно быть достигнуто соглашение (статья~432 ГК~РФ)}
}
\newglossaryentry{term-rag}{
    name={Retrieval-Augmented Generation},
    description={метод генерации ответа большой языковой моделью, при котором перед генерацией выполняется поиск релевантных фрагментов во внешней памяти, а найденные фрагменты добавляются в контекст модели}
}
\newglossaryentry{term-vectordb}{
    name={Векторная база данных},
    description={специализированная система хранения, обеспечивающая эффективный поиск ближайших соседей в пространстве плотных векторных представлений объектов по заданной метрике сходства}
}
\newglossaryentry{term-notconcluded}{
    name={Незаключённый договор},
    description={договор, в отношении которого не достигнуто соглашение по всем существенным условиям либо не соблюдены требования к форме; не порождает гражданско-правовых обязательств между сторонами}
}
\newglossaryentry{term-hallucination}{
    name={Правовая галлюцинация},
    description={ответ большой языковой модели на правовой вопрос, содержащий несуществующие нормы, прецеденты или неверно приписанное содержание реально существующих норм при отсутствии видимых признаков недостоверности}
}
\newglossaryentry{term-chunk}{
    name={Чанк},
    description={смысловой фрагмент исходного документа, получаемый при предварительной обработке корпуса для целей retrieval-augmented generation; типичной единицей чанка служит абзац, раздел или иной осмысленный текстовый блок}
}
\newglossaryentry{term-checkpoint}{
    name={Контрольный пункт (checkpoint)},
    description={сохранённое состояние графа агента в промежуточной точке его прохождения, позволяющее возобновить обработку после прерывания}
}
\newglossaryentry{term-prompt}{
    name={Промпт},
    description={входной текст для большой языковой модели, задающий контекст, инструкции и предъявляемые к ответу требования}
}
\newglossaryentry{term-embedding}{
    name={Эмбеддинг},
    description={плотное векторное представление текстового фрагмента фиксированной размерности, получаемое нейросетевой моделью и сохраняющее семантическое сходство исходных текстов в виде близости соответствующих векторов}
}
\newglossaryentry{term-biencoder}{
    name={Бикодер (bi-encoder)},
    description={модель векторных представлений, при использовании которой запрос и документ кодируются независимо в векторы фиксированной размерности, а семантическая близость измеряется метрикой между этими векторами; обеспечивает высокую скорость поиска ближайших соседей за счёт предвычисления векторов документов}
}
\newglossaryentry{term-crossencoder}{
    name={Кросс-энкодер (cross-encoder)},
    description={модель, принимающая на вход пару (запрос, документ) одновременно и возвращающая скаляр-оценку их релевантности; не масштабируется на полнокорпусный поиск, поэтому применяется на этапе переранжирования отобранных бикодером кандидатов}
}
\newglossaryentry{term-reranker}{
    name={Реранкер (reranker)},
    description={модель второго этапа в двухступенчатом поиске, переоценивающая релевантность кандидатов, полученных бикодером, и пересортировывающая их по уточнённому скору; обычно реализуется как кросс-энкодер}
}
\newglossaryentry{term-refgraph}{
    name={Расширение по графу ссылок},
    description={метод обогащения результатов поиска, при котором на этапе предобработки корпуса из текстов фрагментов извлекаются упоминаемые ими ссылки на другие фрагменты, а на этапе поиска по этим ссылкам в результат подгружаются связанные фрагменты с ограничением по числу переходов}
}
\newglossaryentry{term-llmjudge}{
    name={LLM-as-a-judge},
    description={паттерн использования большой языковой модели в качестве оценщика других ответов или промежуточных артефактов; применяется к задачам, для которых сложно сформулировать формальную метрику}
}
\newglossaryentry{term-jsonschema}{
    name={Структурированный вывод},
    description={режим работы языковой модели, при котором ответ ограничен переданной {JSON~Schema}; обеспечивает гарантированно валидный машинно-читаемый ответ без необходимости пост-обработки и часто применяется для классификаторов и оценщиков на базе языковых моделей}
}
\newglossaryentry{term-vagueref}{
    name={Бланкетная отсылка},
    description={отсылка одной правовой нормы к другому нормативному акту без указания его конкретной статьи или пункта (например, <<в порядке, установленном законом о регистрации>>)}
}
\newglossaryentry{term-idempotence}{
    name={Идемпотентность платежа},
    description={свойство операции инициирования оплаты, при котором повторный вызов с тем же ключом идемпотентности не создаёт дубликата платежа, а возвращает результат ранее выполненной операции}
}
\newglossaryentry{term-strictjudge}{
    name={Узел строгой LLM-фильтрации},
    description={консервативный режим LLM-фильтрации релевантности, действующий по принципу <<при сомнении~--- отбросить>>; применяется для подавления шума на ранних этапах многошагового поиска}
}
